\section{Conclusão}
\label{sec:conclusion}

Este trabalho apresentou um benchmark multi-sensor para preenchimento de
lacunas, avaliando 15 métodos clássicos sob múltiplos níveis de ruído e
relacionando desempenho a uma medida local de complexidade baseada em entropia
de Shannon~\cite{shannon1948mathematical} em diferentes escalas. Como referência
contemporânea, foram incluídos baselines de aprendizado
profundo~\cite{ronneberger2015unet,goodfellow2014generative,kingma2014auto,dosovitskiy2021image}
no sensor Sentinel-2, sob o mesmo protocolo de métricas calculadas estritamente
em pixels ausentes.

Os resultados mostram que a associação entre entropia local e qualidade de
reconstrução depende do método e da janela adotada, o que favorece análises
multi-escala e desencoraja interpretações como regra monotônica universal.
Além disso, a estratificação por sensor e a análise de custo computacional
reforçam que decisões metodológicas devem considerar, simultaneamente, o
contexto de aquisição e restrições práticas.

Como continuidade natural, destacam-se: (i) ampliar a cobertura amostral para
métodos clássicos sob diferentes padrões de lacuna; (ii) incorporar modelos
temporais~\cite{beckers2003eof,alvera2005reconstruction} quando séries estejam
disponíveis; e (iii) investigar estratégias de treinamento e seleção de modelos
de \gls{DL} que utilizem entropia como sinal auxiliar, respeitando dependência
de escala e variações entre sensores.
